\documentclass[10pt,a4paper]{amsart}
\usepackage[margin=19mm]{geometry}
\usepackage{amsmath,amssymb,mathtools}
\usepackage[scr=rsfs,cal=euler]{mathalfa}
\usepackage{xfrac}
\usepackage[unicode,hidelinks,bookmarks=false]{hyperref}
\newcommand{\ie}{i.e.\ }
\newcommand{\cf}{cf.\ }
\newcommand{\resp}{respectively\ }
\newcommand{\Subsection}{\subsection}
\providecommand{\nobreakdash}{\nobreak\hbox{-}}
\setlength{\emergencystretch}{2em}
\multlinegap0pt
\usepackage{amssymb}
\usepackage{enumerate}
\usepackage[shortlabels]{enumitem}
\usepackage{float}
\newtheorem{theorem}{Theorem}
\newtheorem{lemma}{Lemma}
\newcommand{\CM}{\textnormal{CM}}
\newcommand{\HH}{\mathcal{H}}
\newcommand{\RM}{\textnormal{RM}}
\newcommand{\RR}{\mathbb{R}}
\newcommand{\lrb}[1]{\left[#1\right]}
\newcommand{\lrcb}[1]{\left\{#1\right\}}
\newcommand{\lrp}[1]{\left(#1\right)}
\title{Equidistribution of CM points and RM curves}
\author{Erick Ross, Hui Xue}
\date{Source: arXiv:2605.28614v1 (CC BY 4.0)}
\begin{document}
\maketitle

\noindent\emph{Source and licence.} Equidistribution of CM points and RM curves. Version arXiv:2605.28614v1 (CC BY 4.0), distributed by arXiv under the Creative Commons Attribution 4.0 International licence, http://creativecommons.org/licenses/by/4.0/, as recorded in the arXiv licence metadata for this exact version and shown on the abstract page https://arxiv.org/abs/2605.28614v1. Full licence notice: \texttt{licenses/arxiv-2605.28614-CC-BY-4.0.txt}.\par\smallskip

\noindent\textit{Editorial scope.} Counted unit: the article's theorem thm:equid-CM-geodesic (source0.tex lines 350--354). The article's complete original proof of it is delivered with it (source0.tex lines 753--839, one \texttt{proof} environment of 87 lines, which the article places away from the statement). No mathematics is written, repaired or re-derived: the statement and the proof are the article's own lines, in the article's own notation, and no retained line is substituted or re-flowed.  Retained uncounted as the source-local context the proof needs: lines 249--257; lines 512--529; lines 606--615; lines 643--645; lines 664--677; lines 1231--1238; lines 1357--1364.  Nothing was omitted from any retained span, so the package carries no omission record.  No retained line contains a citation, so the wrapper carries no bibliography and no bibliography entry.  No retained line contains a figure environment, an image-inclusion command or drawing code, so no image is delivered and the package declares no asset dependency.  Editorial formatting only, in addition: an AMS \texttt{amsart} preamble that restates the article's own declarations and macros used by the retained lines, namely the environments \texttt{theorem}, \texttt{lemma} on separate counters, together with the standard packages those lines need.  The retained environments print in this document as Theorem 1, Lemma 1 in order of appearance, whereas the article prints the same items with the numbering of its own full document.  The complete native source of the article as distributed in the arXiv e-print for this exact version is bundled unchanged as \texttt{sources/201550/source0.tex}.
\begin{align}
    G_{(A,B,C)}
    &= \{z \in \HH ~:~ A|z|^2+B\,\mathrm{Re}(z) + C = 0 \} \label{eqn:GABC-set-formula}\\
    &= 
    \begin{cases}
        \{\lrp{\frac{-C}{B}}+iy ~:~ y \in (0,\infty)\} & \text{if } A=0, \\
        \{\lrp{\frac{-B}{2A}} + \lrp{\frac{\sqrt{B^2-4AC}}{2A}}e^{i\theta} ~:~ \theta \in (0,\pi)\} & \text{if } A>0.
    \end{cases} 
\end{align}

\begin{lemma} \label{lem:CM-RM-along-G-diophantine} ~
\begin{enumerate}
\item 
Fix a rational geodesic $G_{(A,B,C)}$. Then $[a,b,c]$ lies on $G_{(A,B,C)}$ iff  
\begin{align} \label{eqn:cond--CM-on-geo}
2aC+2cA=bB.
\end{align}
\item Fix a rational geodesic $G_{(A,B,C)}$. Then $\langle a,b,c \rangle$ has  perpendicular intersection with $G_{(A,B,C)}$ iff
\begin{align} \label{eq:cond--RM-on-geo}
2aC+2cA=bB.
\end{align}
\item 
Fix a CM point $[A,B,C]$. Then $\langle a,b,c \rangle$ passes through $[A,B,C]$ iff 
\begin{align} \label{eqn:cond--RM-on-CM}
2aC+2cA=bB.
\end{align}
\end{enumerate}
\end{lemma}

\begin{align}
    \CM^{G}_{\Delta} 
    %
    &= \Big\{  [nQ,nP,-m]  ~:~ n \ge 1,~ \gcd(m,n)=1,~  0 > n^2P^2 + 4Qnm \ge -\Delta
    \Big\} \\
    %
    &= \Big\{  [nQ,nP,-m]  ~:~ n \ge 1,~ \gcd(m,n)=1,~  0 > B mn + C n^2 \ge -\Delta
    \Big\}, \label{eqn:CM-G0BC-set-formula} \\
    \text{where}&\quad  B := 4Q, \qquad C := P^2,
\end{align}

\begin{align} \label{eqn:def-PQR}
    (P,Q,R) := \tfrac{1}{\gcd(D_0,2)} (-2C_0,-B_0,2A_0) 
\end{align}

\begin{align}
    \CM^{G}_{\Delta} 
    %
    &= \Big\{  [nS,~nPb_0+m\tfrac{R}{S},~nPc_0-m\tfrac{Q}{S}]  ~:~ n \ge 1,~ \gcd(m,n)=1,\\
    &\hspace{23mm}  0 >(nPb_0+m\tfrac{R}{S})^2-4(nS)(nPc_0-m\tfrac{Q}{S}) \ge -\Delta
    \Big\} \\
    %
    &= \Big\{[nS,~nPb_0+m\tfrac{R}{S},~nPc_0-m\tfrac{Q}{S}] ~:~ n\ge 1,~ \gcd(m,n)=1,~ \\
    &\hspace{50mm} 0 > A m^2 + B mn + C n^2  \ge -\Delta
    \Big\} \label{eqn:CM-GABC-set-formula} \\
    %
    \text{where}\quad  A&:=\tfrac{R^2}{S^2}, \qquad B := 2Pb_0 \tfrac{R}{S}+4Q, \qquad C := P^2b_0^2-4SPc_0,  \label{eqn:def-ABC-for-CM^G-set-def}\\
     D &:= B^2 -4A C = 16Q^2+16PR = \frac{16}{\gcd(D_0,2)^2} D_0 > 0.
\end{align}

\begin{theorem} 
    \label{thm:equid-W--A=0,B>0}
    Fix $(B,C) \in \RR^2$ with $B > 0$. Then
    \begin{align}
        W_{\Delta}^{(0,B,C)} = \lrcb{\frac{m}{n}  ~:~ n\ge1,~ \gcd(m,n)=1,~ 0 < B mn + C n^2 \le \Delta} 
    \end{align} 
    is equidistributed over $(\frac{-C}{B},\infty)$ with respect to the measure $d\mu = \frac{1}{Bt+C} \,dt$.
\end{theorem}

\begin{theorem} \label{thm:equid-W--A>0,D>0}
Fix $(A,B,C) \in \RR^3$ with $A > 0$ and $D := B^2-4AC > 0$. Then
\begin{align}
    W_{\Delta}^{(A,B,C)} := \Big\{\frac{m}{n} ~:~ n\ge 1,~ \gcd(m,n)=1,~   0 < A m^2 + B mn + C n^2  \le \Delta
    \Big\}
\end{align}
is equidistributed over $(\frac{-B+\sqrt{D}}{2A},\infty] \cup [-\infty,\frac{-B-\sqrt{D}}{2A})$ with respect to the measure $d\mu = \frac{1}{At^2+Bt+C} \,dt$. Here, the interval $(\frac{-B+\sqrt{D}}{2A},\infty] \cup [-\infty,\frac{-B-\sqrt{D}}{2A})$ is understood to wrap around from $\infty$ to $-\infty$.
\end{theorem}

\begin{theorem} \label{thm:equid-CM-geodesic}
    Fix a rational geodesic $G$, and let 
    $\CM^{G}_{\Delta}$ denote the set of CM points along $G$ with discriminant $|D| \le \Delta$. 
    Then as $\Delta \to \infty$, $\CM^{G}_{\Delta}$ is equidistributed along $G$ with respect to the hyperbolic metric.
\end{theorem}

\begin{proof}  Let $\RM^{\cap G}_{\Delta}$ denote the set of intersection points of the RM curves in $\RM^{G}_{\Delta}$ with $G$. We want to show that as $\Delta\to\infty$, $\RM^{\cap G}_{\Delta}$ is equidistributed along $G$.

\noindent\textbf{Case 1: $G$ is a half-line rational geodesic.}

Write $G$ as $G_{(0,B_0,C_0)}$. Then by an identical argument as for \eqref{eqn:CM-G0BC-set-formula} in Theorem \ref{thm:equid-CM-geodesic} (and using the same parameters), we have that
\begin{align}
    \RM^{G}_{\Delta} 
    %
    &= \Big\{  \langle nQ,nP,-m \rangle  ~:~ n \ge 1,~ \gcd(m,n)=1,~  0 < B mn + C n^2 \le \Delta
    \Big\},
\end{align}
and so
\begin{align}
    \mathrm{Im}\,\RM^{\cap G}_{\Delta} 
    %
    &= \Big\{  \frac{\sqrt{4Qnm +P^2n^2}}{2Qn}  ~:~ n \ge 1,~ \gcd(m,n)=1,~  0 < B mn + C n^2 \le \Delta
    \Big\} \\
    &= \Big\{  \sqrt{\frac{4}{B} \lrp{\frac mn} + \frac{4C}{B^2}}  ~:~ n \ge 1,~ \gcd(m,n)=1,~  0 < B mn + C n^2 \le \Delta
    \Big\} \\
    &=  \Big\{  y\lrp{\frac mn}  ~:~ n \ge 1,~ \gcd(m,n)=1,~  0 < B mn + C n^2 \le \Delta
    \Big\} \\
    &= y(W_{\Delta}^{(0,B,C)}), \\
    \text{where }\ y(t) 
    &:= \sqrt{\frac{4}{B} t + \frac{4C}{B^2}}.
\end{align}

Finally, we know from Theorem \ref{thm:equid-W--A=0,B>0} that $W_{\Delta}^{(0,B,C)}$ is $d\mu = \frac{1}{Bt+C}$ equidistributed over $(\frac{-C}{B},\infty)$. And under the reparametrization 
$y(t) = \sqrt{\frac{4}{B} t + \frac{4C}{B^2}}$, this measure $d\mu = \frac{1}{Bt+C}$ over $t \in (\frac{-C}{B},\infty)$ becomes $d\mu = \frac{2}{B} \frac{1}{y} dy$ over $y \in (0,\infty)$. Thus we have that $\mathrm{Im}\,\RM^{\cap G}_{\Delta} = y(W_{\Delta}^{(0,B,C)})$ is equidistributed over $(0,\infty)$ with respect to the hyperbolic metric $\frac{1}{y} dy$, completing the proof.

\noindent\textbf{Case 2: $G$ is a semicircle rational geodesic.}

Write $G$ as $G_{(A_0,B_0,C_0)}$. Then by an identical argument as for \eqref{eqn:CM-GABC-set-formula} in Theorem \ref{thm:equid-CM-geodesic} (and using the same parameters), we have that
\begin{align}
    \RM^{G}_{\Delta} 
    %
    &= \Big\{\langle nS,~nPb_0+m\tfrac{R}{S},~nPc_0-m\tfrac{Q}{S} \rangle ~:~ n\ge 1,~ \gcd(m,n)=1,~ \\
    &\hspace{61mm} 0 < A m^2 + B mn + C n^2  \le \Delta
    \Big\}, \\
    \text{where}\quad  A&:=\tfrac{R^2}{S^2}, \qquad B := 2Pb_0 \tfrac{R}{S}+4Q, \qquad C := P^2b_0^2-4SPc_0,  \\
     D &:= B^2 -4A C = 16Q^2+16PR = \frac{16}{\gcd(D_0,2)^2} D_0 > 0.
\end{align}

Now, let $z$ denote the perpendicular-intersection point of $\langle a,b,c\rangle$ and $G=G_{(A_0,B_0,C_0)}$. By \eqref{eqn:GABC-set-formula}, we know that $z$ satisfies
\begin{align}
    A_0|z|^2 + B_0 \, \mathrm{Re}(z) + C_0 = 0 \qquad \text{and}\qquad a|z|^2 + b \, \mathrm{Re}(z) + c = 0.
\end{align}
Thus
\begin{align}
    &\lrp{\frac{B_0}{A_0} - \frac ba} \mathrm{Re}(z) + \lrp{\frac{C_0}{A_0} - \frac ca} = 0, \\
    \text{and so }\ \ \mathrm{Re}(z) &= \frac{\frac ca -\frac{C_0}{A_0}}{\frac{B_0}{A_0} - \frac ba} = \frac{A_0c-C_0a}{B_0a-A_0b} = \frac{B_0b-4C_0a}{2B_0a-2A_0b} \qquad\quad \text{(by Lemma \ref{lem:CM-RM-along-G-diophantine})} \\
    &= \frac{-Qb+2Pa}{-2Qa-Rb} = \frac{\frac{Q}{R}\lrp{-2Qa-Rb} + 2\frac{Q^2}{R}a + 2Pa}{-2Qa-Rb} \quad\quad \text{(by \eqref{eqn:def-PQR})}\\
    &= \frac{Q}{R} -  \frac{\frac{2a}{R}(Q^2+PR)}{2Qa+Rb}.
\end{align}
Hence for an angle $\theta=\arg_q \langle nS,~nPb_0+m\tfrac{R}{S},~nPc_0-m\tfrac{Q}{S} \rangle \cap G$ in the set $\mathrm{arg}_q\, \RM^{\cap G}_{\Delta}$, we have that 
\begin{align}
    q+r \cos\theta = \frac{Q}{R} -  \frac{\frac{2nS}{R}(Q^2+PR)}{2QnS+R\lrp{nPb_0+m\tfrac{R}{S}}} = 
    \frac{Q}{R} - \frac{\frac{2S}{R}(Q^2+PR)}{\frac{R^2}{S}\lrp{\frac mn} +2QS+RPb_0  },
\end{align}
and so
\begin{align}
    \arg_q \RM^{\cap G}_{\Delta} 
    &= \Big\{\theta \lrp{\frac{m}{n}} ~:~ n\ge 1,~ \gcd(m,n)=1,~   0  < A \lrp{\frac{m}{n}}^2 + B \lrp{\frac{m}{n}} + C  \le \frac{\Delta}{n^2}
    \Big\} \\
    &= \theta(W_{\Delta}^{(A,B,C)}), \\
    \text{where }\ \theta(t) &:= \arccos\lrp{ \frac 1 r \lrb{
        \frac{Q}{R} - \frac{\frac{2S}{R}(Q^2+PR)}{\frac{R^2}{S}t +2QS+RPb_0  }
    - q}}.
\end{align}

Also note that
\begin{align} 
    q = \frac{Q}{R}, &\qquad r^2-q^2 = \frac{P}{R}, \qquad r=\frac{\sqrt{Q^2+PR}}{R},  \\
    \text{so }\ 
    \theta(t) 
    &= \arccos\lrp{ \frac{R}{\sqrt{Q^2+PR}} \lrb{
        \frac{Q}{R} - \frac{\frac{2S}{R}(Q^2+PR)}{\frac{R^2}{S}t +2QS+RPb_0  }
    - \frac QR }} \\
    &= \arccos\lrp{ \frac{1}{4\sqrt{Q^2+PR}} \lrb{
         - \frac{16(Q^2+PR)}{2\frac{R^2}{S^2}t +4Q+\frac{R}{S}Pb_0  }
    }} \\
    &= \arccos\lrp{ \frac{-\sqrt D}{ 2A t+B } }.
\end{align}

Finally, we know from Theorem \ref{thm:equid-W--A>0,D>0} that $W_{\Delta}^{(A,B,C)}$ is $d\mu = \frac{1}{At^2+Bt+C} dt$ equidistributed over  $(\frac{-B+\sqrt{D}}{2A},\infty] \cup [-\infty,\frac{-B-\sqrt{D}}{2A})$. And under the reparametrization 
$\theta(t) = \arccos\lrp{ \frac{-\sqrt D}{ 2A t+B } }$, this measure
$d\mu = \frac{1}{At^2+Bt+C} dt$ over $t \in (\frac{-B+\sqrt{D}}{2A},\infty] \cup [-\infty,\frac{-B-\sqrt{D}}{2A})$ becomes $\frac{2}{\sqrt{D}} \frac{1}{\sin \theta} d\theta$ over $\theta \in (0,\pi)$. Thus we have that $\arg_q \RM^{\cap G}_{\Delta} = \theta(W_{\Delta}^{(A,B,C)})$ is equidistributed over $(0,\pi)$ with respect to the hyperbolic metric $\frac{1}{\sin \theta} d\theta$, completing the proof.
\end{proof}

\end{document}
